\section{پایلوت خطای پیش‌بینی: پرسشی تازه از همین داده}
\label{sec:pe-pilot}

خط \lr{EEG} به \lr{EMG} در \lr{v3} به نتیجه‌ی منفی رسید (فصل~\ref{sec:v3-results}). پیش از تصمیم درباره‌ی مقاله، یک پرسش دیگر از همین داده آزموده شد که تکلیف اصلاً برای آن طراحی شده بود: وقتی جسم سنگین‌تر یا سبک‌تر از انتظار است، \lr{EEG} چه نشان می‌دهد؟ این پایلوت هیچ مدلی آموزش نمی‌دهد. همه‌ی تحلیل‌ها مستقیم روی فایل‌های خام \lr{WS} و جدول \lr{AllLifts} انجام می‌شوند (پوشه‌ی \lr{pe\_pilot} و دفترچه‌ی \lr{pe\_pilot.ipynb} در شاخه‌ی \lr{pilot-prediction-error}).

\subsection{چرا این پرسش}

در سری‌های وزن، وزن جسم بدون خبر بین ۱۶۵، ۳۳۰ و ۶۶۰ گرم عوض می‌شود. فرد نیروی گرفتن را بر اساس وزن دفعه‌ی قبل آماده می‌کند، پس وقتی وزن عوض شده، پیش‌بینی‌اش غلط درمی‌آید. هر نفر ۳۸ بلند کردن سنگین‌تر و ۳۸ بلند کردن سبک‌تر از دفعه‌ی قبل دارد. مقاله‌ی اصلی داده تغییر \lr{EEG} پس از وزن سنگین‌تر از انتظار را فقط برای یک نفر (\lr{P3}) و بدون آزمون آماری نشان داده است. در جستجوی منابع کاری پیدا نشد که این اثر را در سطح گروه و با کنترل حرکت سنجیده باشد.

\subsection{طراحی}

هر بلند کردن یکی از سه برچسب را می‌گیرد:
\begin{itemize}
\item \textbf{مطابق انتظار:} وزن همان وزن دفعه‌ی قبل است.
\item \textbf{سنگین‌تر:} وزن از دفعه‌ی قبل بیشتر است، مثلاً ۶۶۰ گرم پس از ۱۶۵ گرم.
\item \textbf{سبک‌تر:} وزن از دفعه‌ی قبل کمتر است.
\end{itemize}
اولین بلند کردن هر سری برچسب نمی‌گیرد، چون انتظاری پشتش نیست. بلند کردن‌هایی که سطح جسم هم عوض شده کنار گذاشته می‌شوند تا خطای اصطکاک با خطای وزن قاطی نشود.

\paragraph{کنترل اصلی.} بلند کردن سنگین‌تر با بلند کردن مطابق انتظارِ \emph{همان وزن} مقایسه می‌شود: ۶۶۰ گرمِ پس از ۱۶۵ در برابر ۶۶۰ گرمِ پس از ۶۶۰. جسم فیزیکی در دو طرف یکی است و فقط انتظار فرق دارد. تفاوت برای هر وزن جدا حساب و بعد میانگین گرفته می‌شود. در همه‌ی آزمون‌های گروهی واحد آماری شرکت‌کننده است.

\paragraph{چهار تحلیل.}
\begin{enumerate}
\item \textbf{رفتار:} آیا نیروها اثر کلاسیک انتظار وزن را نشان می‌دهند؟ این آزمون درستی برچسب‌ها را می‌سنجد.
\item \textbf{پاسخ خطا:} تفاوت \lr{EEG} (پتانسیل کند، توان آلفا و بتا)، \lr{EMG} و نیرو از ۰٫۵ ثانیه پیش تا ۱٫۵ ثانیه پس از لمس جسم با هر دو انگشت، با \lr{cluster permutation test} روی ۱۲ نفر. \lr{EEG} و حرکت در یک محور زمان آزموده می‌شوند تا معلوم شود کدام زودتر فرق می‌کند.
\item \textbf{انتظار پیش از تماس:} آیا از \lr{EEG} در ۰٫۶ ثانیه‌ی پیش از تماس، وقتی هنوز هیچ تفاوت فیزیکی‌ای نیست، می‌شود وزن دفعه‌ی قبل (۱۶۵ در برابر ۶۶۰ گرم) را رمزگشایی کرد؟ رمزگشا \lr{logistic regression} با \lr{leave-one-series-out} است. همین رمزگشا روی \lr{EMG}، سرعت دست، کانال‌های پیرامونی و \lr{EEG} پیش از \lr{LED} هم کنترل است.
\item \textbf{یادگیری (اکتشافی):} آیا بزرگی پاسخ \lr{EEG} به خطا، اندازه‌ی اصلاح نرخ نیرو در بلند کردن بعدی را پیش‌بینی می‌کند؟
\end{enumerate}

شرط‌های قبول هر تحلیل پیش از دیدن نتیجه‌ها در دفترچه ثبت شد (جدول~\ref{tab:pe-verdict}).

\paragraph{یک نکته درباره‌ی داده.} در این داده دامنه‌ی \lr{EEG} بزرگ است: میانه‌ی \lr{peak-to-peak} در دو ثانیه از کانال‌های مرکزی حدود ۳۰۰ میکروولت است. با آستانه‌ی ثابت ۱۵۰ میکروولت همه‌ی \lr{trial}ها رد می‌شدند. برای همین حذف \lr{artifact} نسبی است: هر بلند کردن که از میانه‌ی همان نفر به‌علاوه‌ی سه انحراف معیار مقاوم بالاتر باشد کنار می‌رود. در هر نفر بین ۳ تا ۲۶ بلند کردن از ۲۹۴ رد شد. پس از حذف، هر نفر ۱۶۹ تا ۱۷۸ بلند کردن مطابق انتظار، ۲۷ تا ۳۴ سنگین‌تر و ۲۷ تا ۳۴ سبک‌تر دارد.

\subsection{نتیجه‌ها}

\paragraph{۱. رفتار: برچسب‌ها درست کار می‌کنند.} وقتی جسم سنگین‌تر از انتظار است، فاز بار در هر ۱۲ نفر بلندتر است (به‌طور میانگین ۱۸۱ میلی‌ثانیه) و نرخ نیروی گرفتن در هر ۱۲ نفر کمتر است. وقتی سبک‌تر است، فاز بار در ۱۱ نفر کوتاه‌تر و نرخ نیرو در ۹ نفر بیشتر است (جدول~\ref{tab:pe-behaviour}). این همان الگوی کلاسیک \lr{Johansson} و \lr{Westling} است: نیرو بر اساس وزن دفعه‌ی قبل برنامه‌ریزی می‌شود.

\begin{table}[H]
\centering
\caption{رفتار: تفاوت بلند کردن پیش‌بینی‌نشده با بلند کردن مطابق انتظار هم‌وزن. ستون جهت مورد انتظار تعداد نفرهایی است که تفاوت‌شان در جهت پیش‌بینی ادبیات است؛ $p$ از آزمون \lr{Wilcoxon} دوطرفه روی ۱۲ نفر.}
\label{tab:pe-behaviour}
\begin{tabular}{llccc}
\toprule
شرط & سنجه & میانگین تفاوت & جهت مورد انتظار & $p$ \\
\midrule
سنگین‌تر & مدت فاز بار & \lr{+۱۸۱} میلی‌ثانیه & ۱۲ از ۱۲ & \lr{۰٫۰۰۰۵} \\
سنگین‌تر & بیشینه‌ی نرخ نیروی گرفتن & \lr{-۷٫۵} نیوتن بر ثانیه & ۱۲ از ۱۲ & \lr{۰٫۰۰۰۵} \\
سنگین‌تر & بیشینه‌ی نیروی گرفتن & \lr{-۰٫۶۱} نیوتن & ۱۱ از ۱۲ & \lr{۰٫۰۲۱} \\
سبک‌تر & مدت فاز بار & \lr{-۵۸} میلی‌ثانیه & ۱۱ از ۱۲ & \lr{۰٫۰۱۶} \\
سبک‌تر & بیشینه‌ی نرخ نیروی گرفتن & \lr{+۲٫۴} نیوتن بر ثانیه & ۹ از ۱۲ & \lr{۰٫۰۱۲} \\
سبک‌تر & بیشینه‌ی نیروی گرفتن & \lr{+۰٫۱۱} نیوتن & ۸ از ۱۲ & \lr{۰٫۵۲} \\
\bottomrule
\end{tabular}

\end{table}

\fig{pe_behaviour.png}{0.95}{رفتار، یک نقطه برای هر نفر. تفاوت پیش‌بینی‌نشده منهای مطابق انتظار برای بیشینه‌ی نرخ نیروی گرفتن، بیشینه‌ی نیروی گرفتن و مدت فاز بار.}

\paragraph{۲. پاسخ خطا: \lr{EEG} پاسخ می‌دهد، ولی بعد از حرکت.} جدول~\ref{tab:pe-error} همه‌ی آزمون‌ها را نشان می‌دهد.
\begin{itemize}
\item \textbf{سبک‌تر از انتظار:} روشن‌ترین اثر \lr{EEG} در این شرط است. ابتدا یک موج منفی روی کانال‌های مرکزی و آهیانه‌ای (\lr{C3}، \lr{Cz}، \lr{CP1}، \lr{Pz}) از ۲۱۰ تا ۳۹۰ میلی‌ثانیه پس از لمس دیده می‌شود ($p$ برابر \lr{۰٫۰۲۷}). بعد یک موج مثبت روی همان ناحیه از ۵۶۰ تا ۸۰۰ میلی‌ثانیه می‌آید ($p$ برابر \lr{۰٫۰۰۳}). توان بتا روی کانال‌های مرکزی چپ از ۴۲۰ میلی‌ثانیه کم می‌شود ($p$ حدود \lr{۰٫۰۴}).
\item \textbf{سنگین‌تر از انتظار:} \lr{EEG} اثر ضعیف‌تری دارد. فقط توان آلفا از ۴۵۰ تا ۸۱۰ میلی‌ثانیه روی ناحیه‌ی وسیعی بالا می‌رود که کانال‌های کنار چشم (\lr{Fp1}، \lr{Fp2}) را هم دربرمی‌گیرد ($p$ برابر \lr{۰٫۰۲۸}).
\item \textbf{ترتیب زمانی:} در هر دو شرط، حرکت زودتر از \lr{EEG} فرق می‌کند. در سنگین‌تر، \lr{EMG} عضله‌ی بازکننده‌ی انگشتان (\lr{ED}) از ۹۰ میلی‌ثانیه \emph{پیش} از لمس و نیروی گرفتن از ۱۱۰ میلی‌ثانیه پس از لمس فرق می‌کند. در سبک‌تر، نیروی گرفتن و بار از حدود ۶۰ میلی‌ثانیه بیشتر است. این \lr{cluster} نیرو در مرز معناداری است ($p$ حدود \lr{۰٫۰۶})، ولی در شکل~\ref{fig:pe_error_lighter.png} روشن دیده می‌شود. لحظه‌ی جدا شدن جسم از میز هم در دو شرط فرق دارد (جدول~\ref{tab:pe-behaviour}).
\end{itemize}

پس \lr{EEG} به خطای وزن پاسخ می‌دهد، ولی هیچ \lr{cluster} \lr{EEG} پیش از تفاوت حرکت شروع نمی‌شود. با این داده نمی‌شود این پاسخ را از پاسخ مغز به خود حرکت اصلاحی جدا کرد. پاسخ‌های حسی به لمس و به جدا شدن جسم از میز هم در دو شرط در زمان‌های متفاوتی رخ می‌دهند.

\paragraph{چند آزمون.} ده آزمون \lr{cluster} انجام شد (دو شرط ضرب در پنج سیگنال) و بین آن‌ها تصحیحی انجام نشد. اگر شش آزمون \lr{EEG} با \lr{Bonferroni} تصحیح شوند (آستانه \lr{۰٫۰۰۸})، فقط موج مثبت سبک‌تر ($p$ برابر \lr{۰٫۰۰۳}) باقی می‌ماند.

\begin{table}[H]
\centering
\caption{پاسخ خطا: \lr{cluster permutation test} روی ۱۲ نفر برای تفاوت پیش‌بینی‌نشده منهای مطابق انتظار هم‌وزن. زمان‌ها نسبت به لمس جسم با هر دو انگشت است. \lr{cluster}های \lr{EEG} روی کانال‌های همسایه و زمان، و \lr{cluster}های \lr{EMG} و نیرو برای هر سیگنال جدا روی زمان ساخته می‌شوند.}
\label{tab:pe-error}
\resizebox{\textwidth}{!}{\begin{tabular}{llccc}
\toprule
شرط & سیگنال & تعداد \lr{cluster} معنادار & شروع اولین \lr{cluster} & کمترین $p$ \\
\midrule
سنگین‌تر & \lr{EEG} کند (\lr{۰٫۵} تا \lr{۳۰} هرتز) & ۰ & & \lr{۰٫۰۷۱} \\
سنگین‌تر & توان آلفا & ۱ & \lr{۴۵۰} میلی‌ثانیه & \lr{۰٫۰۲۸} \\
سنگین‌تر & توان بتا & ۰ & & \lr{۰٫۲۵} \\
سنگین‌تر & \lr{EMG} & ۲ & \lr{-۹۰} میلی‌ثانیه & \lr{۰٫۰۰۲} \\
سنگین‌تر & نیرو & ۲ & \lr{۱۱۰} میلی‌ثانیه & \lr{۰٫۰۰۱} \\
\midrule
سبک‌تر & \lr{EEG} کند (\lr{۰٫۵} تا \lr{۳۰} هرتز) & ۳ & \lr{۲۱۰} میلی‌ثانیه & \lr{۰٫۰۰۳} \\
سبک‌تر & توان آلفا & ۱ & \lr{۷۲۰} میلی‌ثانیه & \lr{۰٫۰۴۴} \\
سبک‌تر & توان بتا & ۲ & \lr{۴۲۰} میلی‌ثانیه & \lr{۰٫۰۴۰} \\
سبک‌تر & \lr{EMG} & ۰ & & \lr{۰٫۱۱} \\
سبک‌تر & نیرو & ۱ & \lr{۷۱۰} میلی‌ثانیه & \lr{۰٫۰۰۷} \\
\bottomrule
\end{tabular}
}
\end{table}

\fig{pe_error_lighter.png}{0.8}{سبک‌تر از انتظار منهای مطابق انتظار، میانگین ۱۲ نفر، لمس در صفر. از بالا: پتانسیل کند، توان آلفا، توان بتا، پوش \lr{EMG} پنج عضله و نیروی گرفتن و بار. نوار خاکستری: \lr{cluster} با $p<$ \lr{۰٫۰۵} (برای \lr{EEG}، \lr{cluster}هایی که \lr{C3} را دربرمی‌گیرند).}

\fig{pe_error_heavier.png}{0.8}{همان شکل برای سنگین‌تر از انتظار.}

\paragraph{۳. انتظار پیش از تماس: رمزگشایی نمی‌شود.} هیچ رمزگشایی وزن دفعه‌ی قبل را بهتر از شانس پیدا نمی‌کند (جدول~\ref{tab:pe-expectation}). \lr{EEG} مرکزی به \lr{AUC} \lr{۰٫۴۷۸} می‌رسد و فقط در ۴ نفر بالای \lr{۰٫۵} است. میانگین \lr{null} این نوع رمزگشا هم کمی زیر \lr{۰٫۵} است (حدود \lr{۰٫۴۷}، یک سوگیری شناخته‌شده‌ی \lr{AUC} روی امتیازهای جمع‌شده از \lr{fold}ها)، ولی \lr{AUC} واقعی از آن هم بالاتر نیست. نکته‌ی جالب این است که در شرط سنگین‌تر، \lr{EMG} عضله‌ی \lr{ED} پیش از لمس فرق دارد (تحلیل ۲). پس انتظار در عضله دیده می‌شود، ولی در \lr{EEG} سطح سر رمزگشایی نمی‌شود.

\begin{table}[H]
\centering
\caption{رمزگشایی وزن دفعه‌ی قبل (۱۶۵ در برابر ۶۶۰ گرم) از ۰٫۶ ثانیه‌ی پیش از تماس، یا از ۰٫۵ ثانیه‌ی پیش از \lr{LED}. $p$ از آزمون \lr{Wilcoxon} یک‌طرفه‌ی \lr{AUC} بالای \lr{۰٫۵} روی ۱۲ نفر.}
\label{tab:pe-expectation}
\begin{tabular}{lccc}
\toprule
ورودی رمزگشا & میانگین \lr{AUC} & نفرهای بالای \lr{۰٫۵} & $p$ \\
\midrule
\lr{EEG} مرکزی، پیش از تماس & \lr{۰٫۴۷۸} & ۴ از ۱۲ & \lr{۰٫۸۳} \\
\lr{EEG} پیرامونی، پیش از تماس & \lr{۰٫۵۱۲} & ۷ از ۱۲ & \lr{۰٫۱۹} \\
\lr{EEG} مرکزی، پیش از \lr{LED} & \lr{۰٫۴۷۴} & ۶ از ۱۲ & \lr{۰٫۷۹} \\
\lr{EMG}، پیش از تماس & \lr{۰٫۵۱۹} & ۷ از ۱۲ & \lr{۰٫۲۳} \\
سرعت دست، پیش از تماس & \lr{۰٫۴۸۵} & ۷ از ۱۲ & \lr{۰٫۶۰} \\
\bottomrule
\end{tabular}

\end{table}

\fig{pe_expectation.png}{0.8}{\lr{AUC} رمزگشایی انتظار، یک نقطه برای هر نفر.}

\paragraph{۴. یادگیری: اثری دیده نشد.} همبستگی جزئی بزرگی پاسخ بتا در ۲۰۰ تا ۶۰۰ میلی‌ثانیه با اندازه‌ی اصلاح نرخ نیرو در بلند کردن بعدی، به‌طور میانگین کمی منفی است (\lr{Fisher z} برابر \lr{-۰٫۰۶۳}). فقط در ۴ نفر از ۱۲ مثبت است ($p$ برابر \lr{۰٫۱۸}). در این تحلیل، نیروی بیشینه، مدت فاز بار و بیشینه‌ی \lr{EMG} همان بلند کردن کنترل شده‌اند. با حدود ۴۰ تا ۴۶ بلند کردن برای هر نفر، این آزمون فقط اثرهای بزرگ را می‌بیند.

\subsection{جمع‌بندی پایلوت}

\begin{table}[H]
\centering
\caption{شرط‌های قبول که پیش از دیدن نتیجه ثبت شده بودند، و نتیجه‌ی هر تحلیل.}
\label{tab:pe-verdict}
\begin{tabular}{p{0.16\textwidth}p{0.42\textwidth}p{0.32\textwidth}}
\toprule
تحلیل & شرط قبول (پیش از دیدن نتیجه) & نتیجه \\
\midrule
۱. رفتار & فاز بار بلندتر و نرخ نیروی کمتر برای سنگین‌تر در بیشتر نفرها & قبول (۱۲ از ۱۲) \\
۲. پاسخ خطا & \lr{cluster} \lr{EEG} روی کانال‌های مرکزی، ترجیحاً پیش از \lr{EMG} و نیرو & نیمه: اثر هست، ولی بعد از تفاوت حرکت شروع می‌شود \\
۳. انتظار پیش از تماس & \lr{AUC} بالای \lr{۰٫۵} با $p<$ \lr{۰٫۰۵} و بهتر از کنترل‌ها & رد \\
۴. یادگیری & فقط اکتشافی & اثری دیده نشد \\
\bottomrule
\end{tabular}

\end{table}

رفتار دقیقاً همان است که ادبیات می‌گوید، پس برچسب‌ها و داده درست کار می‌کنند. \lr{EEG} به خطای وزن پاسخ می‌دهد، به‌ویژه وقتی جسم سبک‌تر از انتظار است. این پاسخ در سطح گروه و با کنترل وزن فیزیکی، نسخه‌ی گسترش‌یافته‌ی نمایش تک‌نفره‌ی مقاله‌ی اصلی داده است. ولی پاسخ \lr{EEG} همیشه بعد از تفاوت حرکت شروع می‌شود و انتظار پیش از تماس هم رمزگشایی نمی‌شود. پس این پایلوت ادعای تازه‌ای نمی‌دهد که از حرکت جدا باشد. با معیارهایی که از پیش تعیین شده بود، این پرسش به‌تنهایی پایه‌ی یک مقاله‌ی قوی نیست.

\newpage
